In this study, we developed DiRLU, a lightweight and scalable reinforcement-learning framework designed to improve IoT security. The framework combines knowledge distillation, which enables a smaller student model to inherit the knowledge of a larger teacher model, with a feature unlearning mechanism that protects sensitive data without requiring retraining. DiRLU can learn efficiently while maintaining privacy, making it ideal for small, resource-constrained IoT devices. Unlike benchmark models limited to 5\% of the BoT-IoT dataset, we utilized a comprehensive 25\% subset, achieving 99.60\% accuracy and a 99.80\% F1 score. These results prove that DiRLU maintains high performance even when working with larger data samples. Our findings show that feature unlearning is both effective and reversible: removing sensitive features reduced the model’s dependence on them while preserving reliability, and restoring these features recovered the original performance. This flexibility ensures that the system can adapt to changing privacy needs without a significant impact on accuracy. The addition of explainable AI further clarified the decision-making process, increasing trust and transparency in model predictions. This helps users and organizations understand why certain decisions are made, which is crucial for security applications. Overall, DiRLU stands out as a practical and adaptable solution for IoT security. It's lightweight design, privacy-preserving capability, and interpretability make it suitable for real-world deployment. Further work will focus on improving its responsiveness to new types of attacks and strengthening defenses against adversarial techniques to support secure and trustworthy IoT solutions.